\problem{2.36}{Centered infinite square well}
For $V=0$ on $-a<x<a$ and $V=\infty$ elsewhere,
the boundary conditions are $\psi(-a)=\psi(a)=0$.
Even solutions have the form $\cos(kx)$ and satisfy
$\cos(ka)=0$, giving $ka=(j+1/2)\pi$.
Odd solutions have the form $\sin(kx)$ and satisfy
$\sin(ka)=0$, giving $ka=j\pi$.
A combined ordering is
\begin{equation}
\boxed{E_n=\frac{n^2\pi^2\hbar^2}{8ma^2}},\qquad
n=1,2,\ldots,
\end{equation}
with
\begin{equation}
\boxed{\psi_n(x)=
\begin{cases}
a^{-1/2}\cos(n\pi x/2a),&n\text{ odd},\\
a^{-1/2}\sin(n\pi x/2a),&n\text{ even}.
\end{cases}}
\end{equation}
Signs can be changed without physical effect.
Each integral $\int_{-a}^a|\psi_n|^2\dd x=1$.
These are the familiar well states for a box of
width $2a$, shifted to the origin. The first three
are an even ground state with no internal node,
an odd first excited state with one internal node,
and an even second excited state with two nodes.

\problem{2.37}{Cubic sine initial wavefunction}
On $0<x<a$, let $\Psi(x,0)=A\sin^3(\pi x/a)$.
Use $\sin^3\theta=(3\sin\theta-\sin3\theta)/4$ and
$\int_0^\pi\sin^6\theta\dd\theta=5\pi/16$.
Normalization yields
\begin{equation}
1=A^2\frac{5a}{16}
\quad\Rightarrow\quad \boxed{A=\frac4{\sqrt{5a}}}.
\end{equation}
The eigenstate decomposition is
\begin{equation}
\Psi(x,0)=\frac3{\sqrt{10}}\phi_1(x)
-\frac1{\sqrt{10}}\phi_3(x).
\end{equation}
Therefore
\begin{equation}
\boxed{\Psi(x,t)=\frac3{\sqrt{10}}\phi_1(x)\ee^{-\ii\Omega t}
-\frac1{\sqrt{10}}\phi_3(x)\ee^{-9\ii\Omega t}}.
\end{equation}
Both $\phi_1,\phi_3$ are even about $x=a/2$.
The density remains symmetric about the midpoint,
so $\boxed{\avg x(t)=a/2}$ for every $t$.
Energy probabilities are $9/10$ and $1/10$:
\begin{equation}
\boxed{\avg H=\frac9{10}E_1+\frac1{10}E_3
=\frac95\hbar\Omega}.
\end{equation}
\qcheck{The density is not stationary, despite a
constant mean position; interference changes its shape.}

\problem{2.38}{Quantum revivals and classical crossing times}
Any normalized well state has expansion
$\Psi(x,t)=\sum_n c_n\phi_n(x)
\ee^{-\ii n^2\Omega t}$.
Choose $T_q=2\pi/\Omega$; then
$\ee^{-\ii n^2\Omega T_q}=\ee^{-2\pi\ii n^2}=1$
for every integer $n$. Hence
\begin{equation}
\boxed{\Psi(x,T_q)=\Psi(x,0)},\qquad
\boxed{T_q=\frac{4ma^2}{\pi\hbar}}.
\end{equation}
A classical particle traverses the box and returns
after a round-trip distance $2a$ at speed
$v=\sqrt{2E/m}$:
\begin{equation}
\boxed{T_{\mathrm{cl}}=\frac{2a}{v}
=2a\sqrt{\frac{m}{2E}}}.
\end{equation}
Equating the two times gives
\begin{equation}
\boxed{E=\frac{\pi^2\hbar^2}{8ma^2}=\frac{E_1}{4}}.
\end{equation}
This is below the lowest allowed energy of the
infinite well. The two concepts of recurrence
are therefore not equal for any stationary
state in the well.

\problem{2.39}{Second derivative of a state with a cusp}
From Problem 2.7,
\begin{equation}
\Psi(x,0)=A
\begin{cases}x,&0<x<a/2,\\a-x,&a/2<x<a,\end{cases}
\quad A=\sqrt{12/a^3}.
\end{equation}
Inside the open interval $(0,a)$, the first
derivative is the step-function expression
\begin{equation}
\boxed{\Psi'(x,0)=A[1-2\theta(x-a/2)]}.
\end{equation}
Using $\theta'=\delta$,
\begin{equation}
\boxed{\Psi''(x,0)=-2A\delta(x-a/2)}.
\end{equation}
This is a distributional derivative, not an ordinary
second derivative at the cusp. The quadratic energy
form gives
\begin{align}
\avg H&=-\frac{\hbar^2}{2m}\int_0^a
\Psi(x,0)\Psi''(x,0)\dd x\\
&=\frac{\hbar^2}{m}A\Psi(a/2,0)
=\frac{\hbar^2}{m}A^2\frac a2
=\boxed{\frac{6\hbar^2}{ma^2}}.
\end{align}
It agrees with Problem 2.7. A rigorous
alternative is the kinetic-energy quadratic form
$\hbar^2\int_0^a|\Psi'|^2/(2m)$.

\problem{2.40}{A quadratic polynomial times a Gaussian}
Set $q=\sqrt{m\omega/\hbar}\,x$.
The initial state is proportional to
$(1-2q)^2\ee^{-q^2/2}$.
Since $H_0=1$, $H_1=2q$, and
$H_2=4q^2-2$,
\begin{equation}
(1-2q)^2=3H_0-2H_1+H_2.
\end{equation}
Using normalized oscillator functions
$\psi_n=\alpha H_n(q)\ee^{-q^2/2}/\sqrt{2^nn!}$
where $\alpha=(m\omega/\pi\hbar)^{1/4}$,
\begin{equation}
\Psi(x,0)=A_{\mathrm{norm}}
\left[3\psi_0-2\sqrt2\psi_1+2\sqrt2\psi_2\right].
\end{equation}
The squared norm of the bracket is
$9+8+8=25$, so $A_{\mathrm{norm}}=1/5$.
The prefactor multiplying the \emph{unnormalized}
polynomial-Gaussian is $A=\alpha/5$.
Thus
\begin{equation}
\boxed{(c_0,c_1,c_2)=
\left(\frac35,-\frac{2\sqrt2}5,
\frac{2\sqrt2}5\right)}.
\end{equation}
Energy probabilities are $9/25$, $8/25$, $8/25$ and
\begin{equation}
\avg H=
\left[\frac9{25}\frac12+\frac8{25}\frac32+
\frac8{25}\frac52\right]\hbar\omega
=\boxed{\frac{73}{50}\hbar\omega}.
\end{equation}
At $t=\pi/\omega$, relative phases are
$\ee^{-\ii n\pi}=(-1)^n$.
The odd $n=1$ coefficient changes sign while
$n=0,2$ do not, so the polynomial becomes
$(1+2q)^2$ up to a global phase.
The smallest positive time is
$\boxed{T=\pi/\omega}$.

\problem{2.41}{The half harmonic oscillator}
For $x>0$, $V=m\omega^2x^2/2$; on $x\le0$
there is an infinite wall. Thus
$\psi(0)=0$ and $\psi(x)$ must decay as
$x\to+\infty$. Extending a half-line solution
oddly to the entire line produces an
odd full-oscillator eigenfunction.
Only $n=1,3,5,\dots$ are permitted.
Reindex $n=2j+1$, $j=0,1,\dots$:
\begin{equation}
\boxed{E_j=\left(2j+\frac32\right)\hbar\omega}.
\end{equation}
The normalized half-line wavefunctions are
$\boxed{\psi_j^{\rm half}(x)=\sqrt2\,
\psi_{2j+1}^{\rm full}(x)}$ for $x>0$.
The factor $\sqrt2$ compensates for
discarding the negative half of a normalized
odd-state density.

\problem{2.42}{Traveling Gaussian packet: envelope and phase speeds}
For $\Psi(x,0)=A\ee^{-ax^2}\ee^{\ii\ell x}$,
normalization still gives $A=(2a/\pi)^{1/4}$.
The initial Fourier distribution is shifted:
$\phi(k)\propto\ee^{-(k-\ell)^2/(4a)}$.
In the free-particle Fourier integral set
$u=k-\ell$ and complete the Gaussian square.
With $\gamma=1+2\ii\hbar at/m$,
\begin{equation}
\boxed{\Psi(x,t)=
\left(\frac{2a}{\pi}\right)^{1/4}
\gamma^{-1/2}
\exp\left[-\frac{a(x-v_gt)^2}{\gamma}\right]
\exp\left[\ii\ell\left(x-\frac{\hbar\ell t}{2m}\right)\right]},
\quad v_g=\frac{\hbar\ell}{m}.
\end{equation}
The density is a translated spreading Gaussian:
\begin{equation}
|\Psi|^2=
\sqrt{\frac{2a}{\pi|\gamma|^2}}\,
\exp\left[-\frac{2a(x-v_gt)^2}{|\gamma|^2}\right].
\end{equation}
Thus the envelope travels at
$v_g=\hbar\ell/m$; its central phase has
$\omega_\ell=\hbar\ell^2/(2m)$, so
the phase speed is
$v_p=\omega_\ell/\ell=\hbar\ell/(2m)=v_g/2$
for $\ell\ne0$. The packet satisfies
$\avg x(t)=v_gt$, $\avg p=\hbar\ell$.
\qcheck{The local phase of the dispersing Gaussian
envelope also evolves, so this $v_p$ refers to
its plane-wave carrier, not to every local
phase contour of the full packet.}

\problem{2.43}{Infinite well with a central delta barrier}
Take $V(x)=\alpha\delta(x)$ for $|x|<a$,
infinite outside, with $\alpha>0$.
Odd states vanish at $x=0$ and satisfy
the delta jump automatically. They remain
\begin{equation}
\psi_o(x)\propto\sin(j\pi x/a),\qquad
\boxed{E_{o,j}=\frac{j^2\pi^2\hbar^2}{2ma^2}}
\quad(j=1,2,\ldots).
\end{equation}
An even state on the right half may be written
$\psi_e(x)=C\sin[k(a-x)]$ for $0<x<a$,
reflected evenly for $x<0$.
At the origin the jump is
$2\psi_e'(0^+)=2m\alpha\psi_e(0)/\hbar^2$,
giving
\begin{equation}
-2Ck\cos(ka)=\frac{2m\alpha}{\hbar^2}C\sin(ka)
\quad\Rightarrow\quad
\boxed{-ka\cot(ka)=\frac{m\alpha a}{\hbar^2}}.
\end{equation}
At $\alpha=0$, the roots are $ka=(j+1/2)\pi$.
At $\alpha\to\infty$, the even roots approach
$ka=j\pi$ from below and become degenerate
with the unaffected odd states. A repulsive
central barrier raises the even-state energies,
consistent with the variational principle.

\problem{2.44}{No degenerate bound levels on a connected line}
Let $\psi_1,\psi_2$ solve the same
one-dimensional stationary equation at energy $E$:
\begin{equation}
-\frac{\hbar^2}{2m}\psi_j''+V\psi_j=E\psi_j.
\end{equation}
Multiply the first by $\psi_2$, the second by
$\psi_1$ and subtract:
$\psi_2\psi_1''-\psi_1\psi_2''=0$.
This is the derivative of the Wronskian
$W=\psi_2\psi_1'-\psi_1\psi_2'$.
Hence $W$ is constant. For ordinary bound states
$\psi_j,\psi_j'\to0$ at an outer boundary, so
$W=0$ everywhere. Wherever $\psi_1\ne0$,
\begin{equation}
\left(\frac{\psi_2}{\psi_1}\right)'=
\frac{\psi_1\psi_2'-\psi_1'\psi_2}{\psi_1^2}=0.
\end{equation}
Thus $\psi_2=C\psi_1$ and continuity extends the
relation through zeros. The solutions are not
linearly independent. Therefore each isolated
bound energy is nondegenerate on a connected
one-dimensional line under these assumptions.

\problem{2.45}{Node interlacing for different bound energies}
Take real eigenfunctions $\psi_m,\psi_n$ with
$E_n>E_m$ and define
\begin{equation}
W(x)=\psi_m'(x)\psi_n(x)-\psi_m(x)\psi_n'(x).
\end{equation}
Use $\psi_j''=(2m/\hbar^2)(V-E_j)\psi_j$
to obtain
\begin{equation}
\boxed{W'(x)=\frac{2m}{\hbar^2}
(E_n-E_m)\psi_m(x)\psi_n(x)}.
\end{equation}
Between successive nodes $x_1<x_2$ of $\psi_m$,
integrate to get
\begin{equation}
\boxed{\psi_m'(x_2)\psi_n(x_2)-
\psi_m'(x_1)\psi_n(x_1)=
\frac{2m(E_n-E_m)}{\hbar^2}
\int_{x_1}^{x_2}\psi_m\psi_n\dd x}.
\end{equation}
Suppose $\psi_n$ has no zero in $(x_1,x_2)$.
Then it has a definite sign, as does $\psi_m$.
If $\psi_m>0$ in this interval, then
$\psi_m'(x_1)>0$ and $\psi_m'(x_2)<0$;
the left side has the \emph{opposite} sign
to $\psi_n$ while the right side has the
\emph{same} sign. This is a contradiction.
The argument with $\psi_m<0$ is analogous.
Thus $\psi_n$ has a node between adjacent
nodes of $\psi_m$. With Sturm--Liouville
boundary conditions this interlacing yields
the increasing node count with energy.
